\documentclass[letterpaper]{article}
\usepackage[letterpaper,margin=1in]{geometry}
\input{wsp/testing/test-case-library.tex}
\begin{document}
\def\TCID{TC-0057}
\def\TCTitle{Interactive error recovery and batch equivalence}
\def\TCPurpose{Verify recoverable syntax, expansion, redirection, resolution, and launch errors discard only one command and retain batch language meaning.}
\def\TCPriority{\TCPriorityCritical}
\def\TCRequirementRef{WSH-REQ-0032, WSH-REQ-0057}
\def\TCDesignRef{ADR-0010, WSH-SPEC-INTERACTIVE-0001 section 7, and WSH-SPEC-LANGUAGE-0001 section 17}
\def\TCFeatureRef{diagnostics, status, command isolation, prompt recovery, fatal output}
\def\TCTestTechnique{Error guessing, state-transition, and mode-comparison testing.}
\def\TCPreconditions{Interactive and redirected sessions use the same tested executable.}
\def\TCEnvironment{Native console plus redirected batch input.}
\def\TCAssumptions{Console output failure is intentionally fatal.}
\def\TCInputData{One case from every recoverable error class followed by a state-observing valid command.}
\def\TCInitialState{A marker variable has a known value.}
\def\TCProcedure{Submit each failure then a valid command, and compare valid-command meaning with batch execution.}
\def\TCExpectedResult{Each diagnostic is visible, no partial effect occurs, the next prompt accepts input, and valid results match batch mode.}
\def\TCPostConditions{No pending source remains.}
\TCTemplate
\end{document}
